\begin{korollar}{Abstände bezüglich der p--Norm}{AbstaendeBezueglichDerPnorm}
Seien $n \in \N$, $p \in [1, \infty]$, $\eps \in \,\! ]0,1[$ und $m \in \N$.\\
Dann gilt:
\begin{enumerate}[(1)]
\item\label{AbstaendeBezueglichDerPnorm1} Ist $p <\infty$ und $m\le\frac 2\pi\eps^2 n^{1-\abs{\frac2p-1}}$,
so gibt es ein injektives $A\in L(\ell_2^m, \ell_p^n)$, sodass gilt
\[\opnorm A
 {\ell_2^m}{\ell_p^n}\opnorm{A\inv}{A\left(\ell_2^m\right)}{\ell_2^m}
\le \frac {1+\eps}{1-\eps}\text.\]
\item\label{AbstaendeBezueglichDerPnorm2} Ist $p = \infty$ und $m \le \eps^2 \log(n)$, 
so gibt es ein injektives $A\in L(\ell_2^m, \ell_\infty^n)$, sodass gilt
\[\opnorm A
 {\ell_2^m}{\ell_\infty^n}\opnorm{A\inv}{A\left(\ell_2^m\right)}{\ell_2^m}
\le \frac {1+\eps}{1-\eps}\text.\]
\end{enumerate}
\end{korollar}
\begin{proof}
(1) Es gelte $p < \infty$ und
$m\le\frac 2\pi\eps^2 n^{1-\abs{\frac2p-1}}$.\\
Ist $p < 2$, so gilt nach \refclaim{AequivalenzkonstantenFuerPNormen}
{AequivalenzkonstantenFuerPNormen3}:
$\norm\cdot_{2,n}\le\norm\cdot_{p,n}\le n^{\frac 1p-\frac 12}\norm\cdot_{2,n}$,
also nach \refclaim{EigenschaftenDesBanachMazurAbstands}
{EigenschaftenDesBanachMazurAbstands3}
\begin{align}
d_{BM}\left(\ell_p^n, \ell_2^n \right) \le n ^{\frac 1 p - \frac 1 2}
= n ^{\abs{\frac 1 p - \frac 1 2}}\text.\label{abdpneq:1}
\end{align}
Ist $p \ge 2$, so gilt nach \refclaim{AequivalenzkonstantenFuerPNormen}
{AequivalenzkonstantenFuerPNormen3}:
$\norm\cdot_{p,n}\le\norm\cdot_{2,n}\le n^{\frac 12-\frac 1p}\norm\cdot_{p,n}$,
also nach \refclaim{EigenschaftenDesBanachMazurAbstands}
{EigenschaftenDesBanachMazurAbstands3}
\begin{align}
d_{BM}\left(\ell_2^n, \ell_p^n\right) \le n ^{\frac 1 2 - \frac 1 p}
= n ^{\abs{\frac 1 2 - \frac 1 p}} = n ^{\abs{\frac 1 p -\frac 12}}\text.
\label{abdpneq:2}
\end{align}
Somit gilt
\[m \le \frac 2 \pi \eps^2 n^{1- \abs{\frac 2 p - 1}}
= \frac 2 \pi \eps^2 \frac n {\left(n^{\abs{\frac 1 p - \frac 1 2}}\right)^2}
\overset{\substack{\eqref{abdpneq:1},\\\eqref{abdpneq:2}}}
\le\frac 2 \pi \eps^2 \frac n {d_{BM}(\ell_p^n, \ell_2^n)^2}\]
und damit sind die Voraussetzungen von
\ref{ExistenzFastIsometrischerAbbildungen} erfüllt.\\
Die Aussage (1) ist genau die Aussage von
\ref{ExistenzFastIsometrischerAbbildungen}.\\
(2) Es gelte $p=\infty$ und $m \le \eps^2 \log(n)$.\\
Seien $(E,\norm\cdot_E):=\ell_\infty^n$, $T:=S^{m-1}$, $x:= (e_i^n)_{i=1}^n$,
\[\bigg((\Omega, \mfa, \Prim), (\gamma,\delta,\widetilde g),(n, m, mn)\bigg)
\text{ ein }(N(0,1), 3)\text{--Tupel zu }(3,(n,m,mn))\] (existiert nach
\ref{ExistenzNormalNTupel}) und $(Z_{\gamma, x}, M_{\delta_T})$ das
Auswertungspaar zu $(\gamma, \delta, x, T)$.\\
Es gilt $x \neq 0$ und wegen $\bprod \cdot x_n = \id_{\R^n}$ gilt
\begin{align}\norm\cdot_E \circ\bprod \cdot x _n
= \norm\cdot_E = \norm\cdot_{\infty, n}
\overset{\refclaim{AequivalenzkonstantenFuerPNormen}
{AequivalenzkonstantenFuerPNormen2}}\le \norm\cdot_{2, n}\text. \tag {$a$}
\end{align}
Ferner ist $\emptyset \neq T$, $T=-T$ und $T \subseteq S^{m-1}$.\hfill $(b)$\\
Schließlich gilt:
\begin{align*}
\E\left(M_{\delta, T} \right) 
&= \E\left(\sup_{t \in S^{m-1}} \bprod t \delta_m \right)
\overset{\ref{IterierteEuklidischeNorm}}= 
\E\left(\norm\cdot_{2,m} \circ \delta \right)
\overset{\ref{ErwartungswertDer2NormEinesStandardnormalverteiltenProzesses}}=
a_m \overset{\refclaim{AbschaetzungEinesGammafunktionsquotienten}
{AbschaetzungEinesGammafunktionsquotienten2}}\le \sqrt m\\
&\overset{\text{Vor.}}\le \eps \sqrt{\log(n)}
\overset{
\refclaim{UntereSchrankeDesErwartungswertsVonAbsolutenN01Ordnungsstatistiken}
{UntereSchrankeDesErwartungswertsVonAbsolutenN01Ordnungsstatistiken2}}\le
\eps\E\left(\gamma_1^\ast \right) \overset{\ref{Ordnungsstatistik}}=
\eps \E\left(\norm\cdot_{\infty, n}\circ\gamma\right)
=\eps\E\left(Z_{\gamma,x}\right)\text. \tag {$c$}
\end{align*}
Wegen $(a), (b), (c)$ sind die Voraussetzungen von
\ref{ExistenzEingeschraenkterFastIsometrien} erfüllt und wegen 
$T = S^{m-1}$ gibt es 
$\emptyset \neq \mca \subseteq L(\ell_2^m, \ell_\infty^n)$ mit
\[\forall\, A \in \mca: A \text{ injektiv } \wedge \, 
\opnorm A
{\ell_2^m}{\ell_\infty^n}\opnorm{A\inv}{A\left(\ell_2^m\right)}{\ell_2^m}
\le \frac {1+\eps}{1-\eps}\text.\]
Wegen $\mca \neq \emptyset$ ist damit insbesondere (2) gezeigt.
\end{proof}